\clearpage
\begin{explanationbox}
\textbf{先看实数：邻边是 $6$ 与 $4$，分段便按 $3:2$ 分配。}
\begin{center}
\begin{tikzpicture}[font=\small]
\draw[black!30,rounded corners=2pt] (0,0) rectangle (14.4,8.2);
\node[anchor=west,font=\small\bfseries] at (0.5,7.72) {角平分线分的是角：$AB:AC=6:4$，所以 $BD:DC=3:2$};
% MODEL 6 4 5
\begin{scope}[shift={(4,2)},scale=1.15]
\coordinate (A) at (4.5,3.9686269666);
\coordinate (B) at (0,0);
\coordinate (C) at (5,0);
\coordinate (D) at (3,0);
\draw[blue,thick] (A)--(B);\draw[teal!80!black,thick] (A)--(C);\draw[thick] (B)--(C);\draw[orange!85!black,very thick] (A)--(D);
\fill (A) circle (1.3pt) node[above=5pt] {$A$};
\fill (B) circle (1.3pt) node[below left=5pt] {$B$};
\fill (C) circle (1.3pt) node[below right=5pt] {$C$};
\fill (D) circle (1.3pt) node[below=5pt] {$D$};
\path (A)--(B) node[midway,above left=4pt] {$AB=6$};
\path (A)--(C) node[midway,above right=4pt] {$AC=4$};
\draw[orange!85!black,thick] (4.0760869565,3.5947708031) arc[start angle=-138.5903778907,delta angle=27.8855668361,radius=0.5652173913];
\draw[orange!85!black,thick] (4.3001654749,3.4399145098) arc[start angle=-110.7048110546,delta angle=27.8855668361,radius=0.5652173913];
\end{scope}
\draw[blue,|-|] ($(B)+(0,-1.0)$)--($(D)+(0,-1.0)$) node[midway,below=3pt] {$BD=3$};
\draw[teal!80!black,|-|] ($(D)+(0,-1.0)$)--($(C)+(0,-1.0)$) node[midway,below=3pt] {$DC=2$};

\end{tikzpicture}
\end{center}

$BC=5$，总共分成 $3+2=5$ 份，所以 $BD=3,DC=2$。
$A$ 处两角相等，不意味着下面两段相等；$AB$ 对应 $BD$，$AC$ 对应 $DC$。

\textbf{同一个顶点出发，三种线各有自己的定义。}
\begin{center}
\begin{tikzpicture}[font=\small]
\draw[black!30,rounded corners=2pt] (0,0) rectangle (14.4,8.2);
\node[anchor=west,font=\small\bfseries] at (0.5,7.72) {一般三角形：角平分线、中线、高分别落在不同点};
% MODEL 6 4 5
\begin{scope}[shift={(4.2,2.1)},scale=1.05]
\coordinate (A) at (4.5,3.9686269666);
\coordinate (B) at (0,0);
\coordinate (C) at (5,0);
\coordinate (D) at (3,0);
\coordinate (M) at (2.5,0);
\coordinate (H) at (4.5,0);
\draw[purple,thick,dashed] (A)--(M);\draw[green!45!black,thick,densely dotted] (A)--(H);
\draw[blue,thick] (A)--(B);\draw[teal!80!black,thick] (A)--(C);\draw[thick] (B)--(C);\draw[orange!85!black,very thick] (A)--(D);
\fill (A) circle (1.3pt) node[above=5pt] {$A$};
\fill (B) circle (1.3pt) node[below left=5pt] {$B$};
\fill (C) circle (1.3pt) node[below right=5pt] {$C$};
\fill (D) circle (1.3pt) node[below=5pt] {$D$};
\pic[draw,angle radius=3mm] {right angle=A--H--B};\node[below left=13pt] at (M) {$M$};\node[below=18pt] at (H) {$H$};
\end{scope}
\node[anchor=west,text=orange!85!black] at (0.8,6.6) {$AD$：角平分线};
\node[anchor=west,text=purple] at (0.8,5.8) {$AM$：中线};
\node[anchor=west,text=green!45!black] at (0.8,5) {$AH$：高};

\end{tikzpicture}
\end{center}

橙线 $AD$ 平分角，紫色虚线 $AM$ 平分边，绿色点线 $AH$ 垂直于 $BC$。
图中 $D,M,H$ 不同，不能仅凭“看上去在中间”就互相替换。

\begin{center}
\begin{tikzpicture}[font=\small]
\draw[black!30,rounded corners=2pt] (0,0) rectangle (14.4,8.2);
\node[anchor=west,font=\small\bfseries] at (0.5,7.72) {等腰时的特殊重合：角平分线也是中线和高};
% MODEL 5 5 6
\begin{scope}[shift={(4.2,2.2)},scale=1]
\coordinate (A) at (3,4);
\coordinate (B) at (0,0);
\coordinate (C) at (6,0);
\coordinate (D) at (3,0);
\coordinate (M) at (3,0);
\coordinate (H) at (3,0);
\draw[blue,thick] (A)--(B);\draw[teal!80!black,thick] (A)--(C);\draw[thick] (B)--(C);\draw[orange!85!black,very thick] (A)--(D);
\fill (A) circle (1.3pt) node[above=5pt] {$A$};
\fill (B) circle (1.3pt) node[below left=5pt] {$B$};
\fill (C) circle (1.3pt) node[below right=5pt] {$C$};
\fill (D) circle (1.3pt) node[below=5pt] {$D=M=H$};
\pic[draw,angle radius=3mm] {right angle=A--D--C};
\path (A)--(B) node[midway,above left=4pt] {$AB=5$};
\path (A)--(C) node[midway,above right=4pt] {$AC=5$};
\draw[orange!85!black,thick] (2.61,3.48) arc[start angle=-126.8698976458,delta angle=36.8698976458,radius=0.65];
\draw[orange!85!black,thick] (3,3.35) arc[start angle=-90,delta angle=36.8698976458,radius=0.65];
\end{scope}

\end{tikzpicture}
\end{center}

两邻边相等时比例变成 $1:1$，角平分线才通过中点，并且垂直于底边。
一般图与等腰图分别说明“通常不同”和“特殊重合”，代替了对三条线位置的长段描述。
\end{explanationbox}
